\documentclass[12pt,a4paper,UTF8]{ctexart}
\usepackage[a4paper,top=22mm,bottom=22mm,left=22mm,right=22mm,headheight=15pt]{geometry}
\newcommand{\Needspace}[1]{\par\begingroup\skip0=\lastskip\vskip-\skip0\dimen0=#1\relax\advance\dimen0 by\skip0\vskip\dimen0\penalty-100\vskip-\dimen0\vskip\skip0\penalty9999\endgroup}
\usepackage{amsmath,amssymb,mathtools,bm}
\usepackage{graphicx,xcolor,booktabs,longtable,array,enumitem,fancyhdr,caption,float,microtype,setspace}
\usepackage[hidelinks,unicode]{hyperref}
\usepackage[most]{tcolorbox}
\definecolor{ContestBlue}{HTML}{17365D}
\hypersetup{pdftitle={第078套},pdfauthor={}}
\setstretch{1.24}
\setlength{\parindent}{2em}
\setlength{\parskip}{0.35em}
\setlength{\emergencystretch}{2em}
\setlist{itemsep=0.35em,topsep=0.4em}
\captionsetup{font=small,labelfont=bf,justification=centering}
\setcounter{secnumdepth}{0}
\pagestyle{fancy}
\fancyhf{}
\fancyhead[L]{\small 第078套}
\fancyfoot[C]{\thepage}
\renewcommand{\headrulewidth}{0.4pt}
\ctexset{
  section={format=\Large\bfseries\color{ContestBlue},beforeskip=1.2em,afterskip=0.7em},
  subsection={format=\large\bfseries\color{ContestBlue},beforeskip=1.0em,afterskip=0.5em},
  subsubsection={format=\normalsize\bfseries,beforeskip=0.8em,afterskip=0.35em}
}
\usepackage{calc}
\newcounter{none}
\let\originalsubsection\subsection
\renewcommand{\subsection}{\Needspace{6\baselineskip}\originalsubsection}
\makeatletter
\newcommand{\tightlist}{\setlength{\itemsep}{0pt}\setlength{\parskip}{0pt}}
\newcommand{\maxwidth}{\ifdim\Gin@nat@width>\linewidth\linewidth\else\Gin@nat@width\fi}
\newcommand{\maxheight}{\ifdim\Gin@nat@height>0.72\textheight 0.72\textheight\else\Gin@nat@height\fi}
\newcommand{\pandocbounded}[1]{\begingroup\setkeys{Gin}{width=\maxwidth,height=\maxheight,keepaspectratio}#1\endgroup}
\makeatother
\begin{document}
\begin{titlepage}
  \centering
  \vspace*{58mm}
  {\Huge\bfseries\color{ContestBlue} 第078套\par}
  \vfill
  \vspace*{28mm}
\end{titlepage}
\section{第 1 题}\label{ux7b2c-1-ux9898}

无限大接地导体平面上有一个半径为 \(R\)
的半球形突起。以半球球心为原点建立直角坐标系 \(Oxyz\)，导体平面为
\(z=0\)。在

\[P(\sqrt2R,0,\sqrt2R)\]

处放置一点电荷 \(Q\)。求该电荷所受静电力在 \(x,y,z\) 三个方向上的分量。

\begin{figure}[H]
\centering
\includegraphics[width=0.72\linewidth,height=\textheight,keepaspectratio,alt={图 1-1　接地平面、半球形突起与点电荷的位置关系。}]{../figures/question-01.pdf}
\caption*{图 1-1　接地平面、半球形突起与点电荷的位置关系。}
\end{figure}

\newpage

\section{第 2 题}\label{ux7b2c-2-ux9898}

半径为 \(r\)、质量为 \(m\) 的均匀薄壁圆筒在水平面上以速度 \(v\)
向右纯滚动。前方有质量同为 \(m\)、厚度为 \(h\)
的长垫块，垫块与地面间光滑，初始静止。圆筒与垫块尖角间摩擦足够大，碰撞完全非弹性；垫块长度足够大，始终不离开水平地面。

\begin{enumerate}
\def\labelenumi{\arabic{enumi}.}
\tightlist
\item
  当 \(h=r/4\) 时，求圆筒能不脱离垫块并翻上垫块的初速度范围。
\item
  求存在某个初速度使圆筒能不脱离垫块并翻上的 \(h\) 的范围。
\end{enumerate}

\begin{figure}[H]
\centering
\includegraphics[width=0.68\linewidth,height=\textheight,keepaspectratio,alt={图 2-1　纯滚动圆筒与可水平移动垫块的碰撞几何。}]{../figures/question-02.pdf}
\caption*{图 2-1　纯滚动圆筒与可水平移动垫块的碰撞几何。}
\end{figure}

\newpage

\section{第 3 题}\label{ux7b2c-3-ux9898}

一根均匀弹簧置于无摩擦水平管中，管子绕一端的竖直轴以恒定角速度
\(\omega\) 转动。弹簧内端固定在距转轴 \(r_0\) 处。弹簧质量为
\(m\)、劲度系数为 \(k\)、自然长度为
\(l_0\)，遵从胡克定律且不会损坏。求稳定伸长后的长度
\(l\)，说明解存在的约束；并讨论 \(r_0\to0\) 的情况。

\begin{figure}[H]
\centering
\includegraphics[width=0.82\linewidth,height=\textheight,keepaspectratio,alt={图 3-1　随径向管匀速转动的有质量弹簧。}]{../figures/question-03.pdf}
\caption*{图 3-1　随径向管匀速转动的有质量弹簧。}
\end{figure}

\newpage

\section{第 4 题}\label{ux7b2c-4-ux9898}

半径 \(r=1\,\mathrm m\) 的均匀电阻圆环固定在水平面内，单位长度电阻
\(\lambda=1\,\Omega/\mathrm m\)。一根单位长度电阻同为 \(\lambda\)
的均匀直棒以圆环直径为长度，其中心与环心重合，可绕过中心的竖直轴无摩擦转动；转动惯量为
\(mr^2\)，\(m=1\,\mathrm{kg}\)。圆环四等分点和转轴均接地。空间有竖直向下的匀强磁场
\(B=1\,\mathrm T\)。初始棒平分相邻接地点，角速度
\(\omega_0=10\,\mathrm{rad/s}\)。

\begin{enumerate}
\def\labelenumi{\arabic{enumi}.}
\tightlist
\item
  求初始流过直棒的电流 \(I_0\)。
\item
  足够长时间后，直棒停在标示的 \(A\) 区还是 \(B\) 区？
\end{enumerate}

\begin{figure}[H]
\centering
\includegraphics[width=0.58\linewidth,height=\textheight,keepaspectratio,alt={图 4-1　接地电阻圆环与可转动导电直棒。}]{../figures/question-04.pdf}
\caption*{图 4-1　接地电阻圆环与可转动导电直棒。}
\end{figure}

\newpage

\section{第 5 题}\label{ux7b2c-5-ux9898}

一线圈密绕成细圆环，圆环半径 \(a=0.1\,\mathrm m\)，线圈截面积
\(S=2\,\mathrm{cm^2}\)，匝数
\(N=200\)，内部为空气。电流在一小段时间内以恒定速率

\[\frac{dI}{dt}=10\,\mathrm{A/s}\]

变化。同时在细圆环的几何中心放置一个初始静止的质子。忽略量子效应与相对论效应，求此时质子的瞬时加速度。

\begin{figure}[H]
\centering
\includegraphics[width=0.84\linewidth,height=\textheight,keepaspectratio,alt={图 5-1　细圆环线圈的空间结构与位于几何中心的质子。}]{../figures/question-05.pdf}
\caption*{图 5-1　细圆环线圈的空间结构与位于几何中心的质子。}
\end{figure}

\newpage

\section{第 6 题}\label{ux7b2c-6-ux9898}

激光陀螺利用沿闭合光路传播的两束反向激光测量相对于惯性空间的转动。先考虑中心到边的距离为
\(r_0\)、周长为 \(L\)、面积为 \(A=r_0L/2\)
的正方形环路，环路绕垂直于平面的轴以角速度 \(\omega\) 逆时针转动。

\begin{enumerate}
\def\labelenumi{\arabic{enumi}.}
\tightlist
\item
  求逆时针光束绕一圈的光程 \(L_+\)。
\item
  求顺、逆两束光的光程差，并在 \(r_0\omega\ll c\) 下保留一阶项。
\item
  对
  \(\omega=15^\circ/\mathrm h\)、\(\lambda=0.70\,\mu\mathrm m\)、\(A=600\times300\,\mathrm{m^2}\)，说明光程差约为
  \(\lambda/4\)。
\item
  改成闭合环形共振腔后，求 \(L_+-L_-\)。
\item
  用 \(A,\lambda,L,\omega\) 表示两束光的拍频 \(f_--f_+\)。
\end{enumerate}

\begin{figure}[H]
\centering
\includegraphics[width=0.523\linewidth,height=\textheight,keepaspectratio,alt={图 6-1　随平台转动的正方形闭合光路。}]{../figures/question-06.pdf}
\caption*{图 6-1　随平台转动的正方形闭合光路。}
\end{figure}

\newpage

\section{第 7 题}\label{ux7b2c-7-ux9898}

把室内空调视为工作在室外温度 \(T_H\) 与室内温度 \(T_L\)
之间的理想可逆制冷机。每循环从室内吸热 \(Q_L\)，向室外放热
\(Q_H\)，电泵做功 \(W\)。空调输入功率为 \(P\)。

\begin{enumerate}
\def\labelenumi{\arabic{enumi}.}
\tightlist
\item
  用 \(T_H,T_L,P\) 表示最大室内热量移除速率。
\item
  若室外以速率 \(\dot Q_{\rm in}=k(T_H-T_L)\)
  向室内传热，求稳态最低室温。
\item
  \(k=173\,\mathrm{W/K}\)、室外 \(40^\circ\mathrm C\)、室内目标
  \(25^\circ\mathrm C\)，求所需最小功率。
\end{enumerate}

\begin{figure}[H]
\centering
\includegraphics[width=0.76\linewidth,height=\textheight,keepaspectratio,alt={图 7-1　空调从室内吸热并向室外放热的能量流。}]{../figures/question-07.pdf}
\caption*{图 7-1　空调从室内吸热并向室外放热的能量流。}
\end{figure}

\newpage

\section{第 8 题}\label{ux7b2c-8-ux9898}

在惯性系 \(S\) 中，一长直圆筒半径为 \(R\)，固有表面电荷面密度为
\(\lambda_1/(2\pi R)\)（按扫描参考解口径；若题给的是 \(S\)
系观测密度，见解析中的另一口径），圆筒以速度 \(v\) 沿 \(+z\)
方向运动；中轴处有线电荷密度为 \(\lambda_2\) 的长直导线，在 \(S\)
中静止。点 \(P\) 距中轴 \(d\)。惯性系 \(S'\) 相对 \(S\) 以速度 \(u\) 沿
\(+z\) 方向运动。

\begin{enumerate}
\def\labelenumi{\arabic{enumi}.}
\tightlist
\item
  在 \(S\) 中分别求 \(d>R\) 与 \(d<R\) 时 \(P\) 点的电场和磁场。
\item
  求 \(S'\) 中圆筒、直导线的速度与线电荷密度。
\item
  当 \(d>R\) 时，求使 \(S'\) 中只有电场而无磁场的 \(u\) 及存在条件。
\item
  求使 \(S'\) 中只有磁场而无电场的 \(u\) 及存在条件。
\item
  是否存在某个 \(S'\) 使电场、磁场同时为零？
\end{enumerate}

\begin{figure}[H]
\centering
\includegraphics[width=0.701\linewidth,height=\textheight,keepaspectratio,alt={图 8-1　运动带电圆筒、静止轴线电荷与共轴运动参考系。}]{../figures/question-08.pdf}
\caption*{图 8-1　运动带电圆筒、静止轴线电荷与共轴运动参考系。}
\end{figure}
\end{document}
